% Independent preview; original geometry and numeric relationships retained.
% Original diagram showing the three components of a directed acyclic graph Markov blanket.
\begin{tikzpicture}[x=1mm,y=1mm,font=\small,
  pgm variable/.style={circle,draw=FigureInk,fill=FigurePaper,
    minimum size=8mm,inner sep=0pt},
  center/.style={pgm variable,draw=FigureRed,fill=FigureRed!18!white,
    very thick},
  blanket/.style={pgm variable,draw=LancetGreen,fill=FigureYellow!10!white,
    thick},
  outside/.style={pgm variable,draw=FigureMuted,fill=FigurePaper,
    dashed},
  edge/.style={-{Stealth[length=1.8mm]},draw=FigureInk,thick},
  outside edge/.style={-{Stealth[length=1.8mm]},draw=FigureMuted,dashed}]
  \node[outside] (o) at (104,56) {$O$};
  \node[blanket] (p1) at (22,50) {$P_1$};
  \node[blanket] (p2) at (74,50) {$P_2$};
  \node[center] (x) at (53,29) {$X_i$};
  \node[blanket] (c1) at (36,7) {$C_1$};
  \node[blanket] (c2) at (76,7) {$C_2$};
  \node[blanket] (u1) at (7,7) {$U_1$};
  \node[blanket] (u2) at (106,7) {$U_2$};

  \draw[outside edge] (o)--(p2);
  \draw[edge] (p1)--(x);
  \draw[edge] (p2)--(x);
  \draw[edge] (x)--(c1);
  \draw[edge] (x)--(c2);
  \draw[edge] (u1)--(c1);
  \draw[edge] (u2)--(c2);

  \node[font=\footnotesize,text=FigureMuted,font=\FigureNoteFont\footnotesize] at (48,61)
    {$P_1,P_2$：父节点};
  \node[font=\footnotesize,text=FigureMuted,font=\FigureNoteFont\footnotesize] at (55,-5)
    {$C_1,C_2$：子节点\qquad $U_1,U_2$：子节点的其他父节点};
  \node[anchor=east,font=\footnotesize,text=FigureMuted,font=\FigureNoteFont\footnotesize] at (99,65)
    {$O$：毯外节点};
\end{tikzpicture}
